% ============================================================
%  Notação Científica — Apostila Premium | PratiqueJá
%  Engine: XeLaTeX
% ============================================================
\documentclass[12pt]{article}
\usepackage{../../pratiqueja}

% \hlm — destaque de resultado em modo-math (cor sem bold)
\newcommand{\hlm}[1]{{\color{darkblue}#1}}

\setsubject{Notação Científica}

\begin{document}
\pagenumbering{arabic}
\setstretch{1.2}
\color{bodytext}

\heroheader{Notação Científica}%
           {Forma $a \times 10^n$, Conversão e Operações}%
           {Matemática · Intermediário}

\setlength{\columnsep}{18pt}
\setlength{\columnseprule}{0.5pt}
\def\columnseprulecolor{\color{exborder}}

\begin{multicols}{2}

%─────────────────────────────────────────────────────────────
\TW{Def}{badgepurple}{Definição}{Forma $a\times10^n$ com $1\leq a<10$}

\regra{A \textbf{notação científica} representa números muito grandes
ou muito pequenos na forma $a \times 10^n$, com $1 \leq a < 10$ e $n$
inteiro.}

\formula{$a \times 10^{n}, \qquad 1 \leq a < 10$}

\exemp{%
  $n > 0$ → grande: $3{,}2 \times 10^6 = 3\,200\,000$\\[3pt]
  $n < 0$ → pequeno: $5 \times 10^{-4} = 0{,}0005$\\[3pt]
  $n = 0$ → $a \times 10^0 = a$%
}

\nota{O expoente $n$ indica quantas casas a vírgula se deslocou da
posição original.}

%─────────────────────────────────────────────────────────────
\TW{$\leftrightarrow$}{darkblue}{Conversão}{Decimal $\leftrightarrow$ notação científica}

\regra{\textbf{Número → científica:} desloque a vírgula até restar
\textbf{1 algarismo} antes dela. O expoente é o nº de casas
(esquerda $\Rightarrow +$, direita $\Rightarrow -$).}

\SH{Número grande}
\exemp{%
  $186\,000\,000 = \hlm{1{,}86 \times 10^8}$ {\footnotesize(8 casas à esq.)}\\[3pt]
  $45\,700 = \hlm{4{,}57 \times 10^4}$%
}

\SH{Número pequeno}
\exemp{%
  $0{,}000032 = \hlm{3{,}2 \times 10^{-5}}$ {\footnotesize(5 casas à dir.)}\\[3pt]
  $0{,}008 = \hlm{8 \times 10^{-3}}$%
}

\SH{Científica → número}
\exemp{%
  $6{,}02 \times 10^{23}$ → 602 seguido de 21 zeros\\[3pt]
  $9{,}1 \times 10^{-31}$ → $0{,}\underbrace{00\ldots0}_{30}91$%
}

%─────────────────────────────────────────────────────────────
\TW{$\times\,\div$}{darkblue}{Multiplicação e Divisão}{Some ou subtraia os expoentes}

\formula{$(a\times10^m)(b\times10^n) = (a\cdot b)\times10^{m+n}$\\[3pt]
$(a\times10^m)\div(b\times10^n) = (a\div b)\times10^{m-n}$}

\obs{Se $a\cdot b$ sair de $[1,10)$, ajuste a vírgula e corrija o
expoente.}

\exemp{%
  \textbf{Multiplicação:}\\
  $(3\times10^4)(2\times10^3) = \hlm{6\times10^7}$\\[3pt]
  \textbf{Com ajuste:}\\
  $(4\times10^5)(3\times10^2) = 12\times10^7 = \hlm{1{,}2\times10^8}$\\[3pt]
  \textbf{Divisão:}\\
  $(9\times10^6)\div(3\times10^2) = \hlm{3\times10^4}$%
}

%─────────────────────────────────────────────────────────────
\TW{$+\,-$}{darkblue}{Adição e Subtração}{Iguale os expoentes antes de operar}

\regra{Os expoentes devem ser \textbf{iguais}. Ajuste um dos termos,
opere e simplifique.}

\exemp{%
  $3{,}2\times10^5 + 4{,}0\times10^4$\\
  $= 3{,}2\times10^5 + 0{,}4\times10^5$ {\footnotesize(igualando)}\\
  $= \hlm{3{,}6\times10^5}$%
}

%─────────────────────────────────────────────────────────────
\needspace{11\baselineskip}
\TW{Aplic.}{badgegreen}{Aplicações}{Ciência, tecnologia e ENEM}

\exemp{%
  \textbf{Distância Terra–Sol:} $1{,}496\times10^{11}\,$m\\[2pt]
  \textbf{Massa do próton:} $1{,}67\times10^{-27}\,$kg\\[2pt]
  \textbf{Tamanho de um vírus:} $10^{-7}$ a $10^{-6}\,$m\\[2pt]
  \textbf{Velocidade da luz:} $3\times10^8\,$m/s%
}

\SH{Problema — ENEM}
\exemp{%
  Uma colônia de bactérias dobra a cada hora. Iniciando com
  $5\times10^3$, após 4 horas:\\[2pt]
  $5\times10^3 \times 2^4 = 5\times10^3 \times 16$\\
  $= 80\times10^3 = \hlm{8\times10^4}$ bactérias%
}

\end{multicols}

\end{document}
